% ============================================================================
%  Section VIII — Conclusion  (IOP MST)
%  Passive voice / third person throughout. Requires: none beyond base.
% ============================================================================

\section{Conclusion}
\label{sec:conclusion}

This paper examined whether low-cost, low-rate smartphone-MEMS sensing can support
condition monitoring of induction machines, and how such monitoring should be
evaluated and deployed. Using a public $100$~Hz vibration dataset, the near-perfect
accuracies reported under random-split evaluation were shown to be largely an
artefact of data leakage: under a leakage-free, recording-wise protocol the
accuracy of standard classifiers fell by approximately $41$ to $47$ percentage
points, with comparable reductions under cross-speed and cross-load shifts. The
residual accuracy nonetheless remained well above chance, confirming that coarse
health states---notably a healthy machine and a broken rotor bar---are recoverable
at $100$~Hz, whereas the fine staging of incipient bearing degradation, whose
signatures lie beyond the $50$~Hz usable band, is not.

Characterising the consumer sensor as a measurement instrument showed that its
on-device gravity compensation preserves the vibration power but reshapes its
low-frequency structure, leaving the raw channels the more informative, so that
they should be preferred. Most consequentially for
deployment, the accuracy--resource analysis established that a resource-frugal
configuration---$25$~Hz sampling, three raw channels, and a $4$~s window---not
only suffices but outperforms the naive configuration under honest evaluation,
raising accuracy to $76\%$ while reducing the raw data volume roughly eight-fold;
combined with a communication link budget and a kilobyte-scale edge-model
footprint, this yields concrete operating points for Bluetooth Low Energy,
Zigbee, NB-IoT, and LoRaWAN nodes. Under honest evaluation, therefore, a
smartphone-class sampling rate is not in itself a barrier to IoT-based machine
monitoring; on this dataset, moderate decimation to $25$~Hz did not reduce---and
in fact improved---the measured accuracy, although whether this reflects a
genuine anti-aliasing benefit or a dataset-specific effect remains to be
established on further machines.

Future work will extend the evaluation across multiple machines of different
ratings and manufacturers to establish cross-machine generalisation, and will
complement the present analysis with direct, cycle-accurate measurement of
on-device inference latency and energy. The evaluation protocol set out in this
paper is intended to serve as a deployment-honest baseline for these efforts.
